% ============================================================
% figures/styles.tex
% Styles communs à toutes les figures vectorielles du rapport.
%
% Règles visuelles, tirées de images_patterns_accepted/readme.txt :
%   - strictement monochrome : noir, blanc, nuances de gris
%   - fond blanc
%   - typographie propre et lisible, sans empattement pour les
%     figures afin de les distinguer du corps du texte
%   - angles nets, alignements réguliers, espacement uniforme
%   - notation UML respectée, chaque type de diagramme gardant
%     son identité propre
% ============================================================

\usetikzlibrary{
  shapes.geometric,
  shapes.multipart,
  shapes.misc,
  arrows.meta,
  positioning,
  calc,
  fit,
  backgrounds,
  patterns,
  decorations.pathreplacing,
  chains,
  matrix
}

% ── Nuances de gris ─────────────────────────────────────────
\definecolor{figGrisTresClair}{gray}{0.94}
\definecolor{figGrisClair}{gray}{0.86}
\definecolor{figGrisMoyen}{gray}{0.70}
\definecolor{figGrisFonce}{gray}{0.45}

% ── Réglages généraux ───────────────────────────────────────
\tikzset{
  every picture/.style={
    line width=0.5pt,
    font=\sffamily\scriptsize
  },
  %
  % ── Boîtes génériques ─────────────────────────────────────
  fbox/.style={
    draw=black, fill=white, rectangle,
    minimum height=7mm, minimum width=22mm,
    inner sep=2pt, align=center
  },
  fboxg/.style={fbox, fill=figGrisTresClair},
  fboxgg/.style={fbox, fill=figGrisClair},
  fboxtitre/.style={fbox, fill=figGrisClair, font=\sffamily\scriptsize\bfseries},
  %
  % ── Zones de regroupement ─────────────────────────────────
  zone/.style={
    draw=figGrisFonce, dashed, rounded corners=2pt,
    inner sep=3mm
  },
  zonepleine/.style={
    draw=black, rounded corners=2pt, inner sep=3mm, fill=figGrisTresClair
  },
  etiquettezone/.style={
    font=\sffamily\scriptsize\itshape, text=black, fill=white, inner sep=1pt
  },
  %
  % ── UML : cas d'utilisation ───────────────────────────────
  cu/.style={
    draw=black, fill=white, ellipse,
    minimum width=30mm, minimum height=9mm,
    inner sep=1pt, align=center
  },
  %
  % ── UML : classes ─────────────────────────────────────────
  classe/.style={
    draw=black, fill=white,
    rectangle split, rectangle split parts=3,
    rectangle split draw splits=true,
    minimum width=28mm, inner sep=2pt,
    text width=28mm, align=left,
    every one node part/.style={font=\sffamily\scriptsize\bfseries, align=center}
  },
  classe2/.style={
    draw=black, fill=white,
    rectangle split, rectangle split parts=2,
    rectangle split draw splits=true,
    minimum width=28mm, inner sep=2pt,
    text width=28mm, align=left,
    every one node part/.style={font=\sffamily\scriptsize\bfseries, align=center}
  },
  classetitre/.style={
    draw=black, fill=figGrisTresClair, rectangle,
    minimum width=28mm, inner sep=3pt, align=center,
    font=\sffamily\scriptsize\bfseries
  },
  interface/.style={
    classe2, fill=figGrisTresClair
  },
  %
  % ── UML : activités ───────────────────────────────────────
  action/.style={
    draw=black, fill=white, rectangle, rounded corners=3pt,
    minimum height=7mm, minimum width=26mm,
    inner sep=2pt, align=center
  },
  decision/.style={
    draw=black, fill=white, diamond, aspect=2.4,
    inner sep=1pt, align=center, minimum height=7mm
  },
  debut/.style={
    draw=black, fill=black, circle, minimum size=3mm, inner sep=0pt
  },
  fin/.style={
    draw=black, fill=white, circle, minimum size=4mm, inner sep=0pt,
    path picture={\fill (path picture bounding box.center) circle (1.1mm);}
  },
  barre/.style={
    draw=black, fill=black, rectangle,
    minimum width=30mm, minimum height=1.2mm, inner sep=0pt
  },
  %
  % ── UML : états ───────────────────────────────────────────
  etat/.style={
    draw=black, fill=white, rectangle, rounded corners=4pt,
    minimum height=8mm, minimum width=26mm,
    inner sep=2pt, align=center
  },
  %
  % ── UML : composants et déploiement ───────────────────────
  composant/.style={
    draw=black, fill=white, rectangle,
    minimum height=9mm, minimum width=30mm,
    inner sep=3pt, align=center
  },
  noeud/.style={
    draw=black, fill=white, rectangle,
    minimum height=10mm, minimum width=30mm,
    inner sep=3pt, align=center
  },
  paquetage/.style={
    draw=black, fill=white, rectangle,
    minimum height=10mm, minimum width=30mm,
    inner sep=3pt, align=center
  },
  %
  % ── UML : séquence ────────────────────────────────────────
  entete/.style={
    draw=black, fill=figGrisTresClair, rectangle,
    minimum height=7mm, minimum width=24mm,
    inner sep=2pt, align=center, font=\sffamily\scriptsize\bfseries
  },
  ligne/.style={draw=black, dashed, dash pattern=on 1.6pt off 1.6pt},
  activation/.style={draw=black, fill=figGrisClair, minimum width=2.4mm, inner sep=0pt},
  %
  % ── Flèches ───────────────────────────────────────────────
  flux/.style={-{Latex[length=2mm, width=1.4mm]}, draw=black},
  fluxep/.style={-{Latex[length=2.4mm, width=1.8mm]}, draw=black, line width=0.9pt},
  assoc/.style={draw=black},
  dirige/.style={-{Straight Barb[length=1.6mm]}, draw=black},
  dependance/.style={-{Straight Barb[length=1.6mm]}, draw=black, dashed,
                     dash pattern=on 1.8pt off 1.8pt},
  heritage/.style={-{Triangle[open, length=2.6mm, width=2.6mm]}, draw=black},
  realisation/.style={-{Triangle[open, length=2.6mm, width=2.6mm]}, draw=black, dashed,
                      dash pattern=on 1.8pt off 1.8pt},
  composition/.style={{Diamond[length=2.6mm, width=1.8mm]}-, draw=black},
  agregation/.style={{Diamond[open, length=2.6mm, width=1.8mm]}-, draw=black},
  message/.style={-{Latex[length=1.8mm, width=1.3mm]}, draw=black},
  retour/.style={-{Latex[length=1.8mm, width=1.3mm]}, draw=black, dashed,
                 dash pattern=on 1.6pt off 1.6pt},
  %
  % ── Étiquettes ────────────────────────────────────────────
  % align est indispensable : sans lui, un \\ dans un libellé de
  % noeud provoque une erreur de compilation.
  etiq/.style={font=\sffamily\tiny, fill=white, inner sep=1pt, align=center},
  etiqi/.style={font=\sffamily\tiny\itshape, fill=white, inner sep=1pt, align=center},
  card/.style={font=\sffamily\tiny, inner sep=1.5pt, align=center},
  note/.style={
    draw=black, fill=figGrisTresClair, rectangle,
    inner sep=3pt, align=left, font=\sffamily\tiny\itshape,
    rounded corners=1pt
  }
}

% ── Acteur UML : bonhomme filaire ───────────────────────────
% Usage : \acteur{nom_du_noeud}{position}{libellé}
\newcommand{\acteur}[3]{%
  \begin{scope}[shift={#2}]
    \draw (0,0.62) circle (0.155);
    \draw (0,0.465) -- (0,0.06);
    \draw (-0.22,0.36) -- (0.22,0.36);
    \draw (0,0.06) -- (-0.19,-0.30);
    \draw (0,0.06) -- (0.19,-0.30);
    \node[font=\sffamily\scriptsize, align=center, anchor=north] at (0,-0.36) {#3};
    \coordinate (#1) at (0,0.2);
    \coordinate (#1-bas) at (0,-0.36);
    \coordinate (#1-droite) at (0.28,0.2);
    \coordinate (#1-gauche) at (-0.28,0.2);
  \end{scope}%
}

% ── Symbole de composant UML (petit rectangle à onglets) ────
\newcommand{\symbcomposant}[1]{%
  \begin{scope}[shift={#1}, scale=0.9]
    \draw (0,0) rectangle (0.34,0.24);
    \draw (-0.07,0.045) rectangle (0.10,0.095);
    \draw (-0.07,0.145) rectangle (0.10,0.195);
  \end{scope}%
}

% ── Onglet de paquetage UML ─────────────────────────────────
% À appeler APRÈS avoir posé le noeud, qui conserve ses ancres.
% Usage : \node[fbox, minimum height=11mm] (p) {...}; \ongletPaquet{p}
\newcommand{\ongletPaquet}[1]{%
  \draw[fill=white]
    ([yshift=-0.2pt]#1.north west) rectangle
    ([xshift=13mm, yshift=3.4mm]#1.north west);
}

% ── Icône de composant, posée dans le coin d'un noeud ───────
\newcommand{\iconeComposant}[1]{%
  \symbcomposant{($(#1.north east)+(-0.58,-0.44)$)}%
}

% ── Interface fournie (boule) et requise (demi-cercle) ──────
% Usage : \interfaceFournie{depuis}{vers} pose la boule au bout.
\newcommand{\interfaceFournie}[2]{%
  \draw (#1) -- (#2);
  \draw[fill=white] (#2) circle (1.15mm);%
}
\newcommand{\interfaceRequise}[3]{%
  % #1 point de départ, #2 point d'accroche, #3 angle d'ouverture
  \draw (#1) -- (#2);
  \draw ([shift={(#3:1.7mm)}]#2) arc (#3:{#3+230}:1.7mm);%
}

% ── Noeud de déploiement UML (boîte en perspective) ─────────
% Usage : \node[noeud] (n) {...}; \perspective{n}
\newcommand{\perspective}[1]{%
  \draw ([shift={(0.22,0.22)}]#1.north west) -- ([shift={(0.22,0.22)}]#1.north east)
        -- ([shift={(0.22,0.22)}]#1.south east) -- (#1.south east);
  \draw (#1.north west) -- ([shift={(0.22,0.22)}]#1.north west);
  \draw (#1.north east) -- ([shift={(0.22,0.22)}]#1.north east);%
}
